\nuevaReceta{Cobbler de Melocotón}{10--12 raciones}{1 h aprox.}{receta:cobbler_melocoton}

\begin{minipage}[t]{0.35\textwidth}

    \section*{Ingredientes}

    \textbf{Relleno}

    \begin{itemize}[leftmargin=*]

        \item 8 melocotones, pelados, deshuesados y cortados en rodajas

        \item 50 g de azúcar blanco

        \item 55 g de azúcar moreno

        \item 1/2 cucharadita de canela molida

        \item 1/2 cucharadita de jengibre molido

        \item 1/4 cucharadita de nuez moscada molida

        \item 2 cucharaditas de maicena

    \end{itemize}

    \textbf{Cobertura}

    \begin{itemize}[leftmargin=*]

        \item 125 g de harina de trigo

        \item 50 g de azúcar blanco

        \item 55 g de azúcar moreno

        \item 1 cucharadita de levadura química (aprox. 4 g)

        \item 1/2 cucharadita de sal

        \item 1/2 cucharadita de canela molida

        \item 1/2 cucharadita de jengibre molido

        \item 113 g de mantequilla fría

        \item 120 ml de agua caliente

    \end{itemize}

\end{minipage}
\hfill
\begin{minipage}[t]{0.6\textwidth}

    \section*{Preparación}

    \begin{enumerate}

        \item \textbf{Precalentar:} Precalentar el horno a 200 °C.

        \item \textbf{El relleno:} Colocar los melocotones en un bol y añadir el azúcar blanco, el azúcar moreno, la canela, el jengibre, la nuez moscada y la maicena.

        \item \textbf{Mezclar:} Remover bien hasta que los melocotones queden cubiertos por el azúcar y las especias. Pasar la mezcla a una fuente de horno de unos 23 x 33 cm.

        \item \textbf{La cobertura:} Mezclar en otro bol la harina, el azúcar blanco, el azúcar moreno, la levadura química, la sal, la canela y el jengibre.

        \item \textbf{La mantequilla:} Cortar la mantequilla fría en dados pequeños y añadirla a la mezcla. Trabajar con las manos hasta obtener una textura similar a migas gruesas.

        \item \textbf{La masa:} Añadir el agua caliente y mezclar hasta formar una masa húmeda.

        \item \textbf{Montaje:} Repartir la masa a cucharadas sobre los melocotones. No hace falta cubrir toda la fruta ni alisar la superficie.

        \item \textbf{Hornear:} Hornear durante 35--40 minutos, hasta que la cobertura esté dorada y el jugo de los melocotones burbujee por los bordes.

        \item \textbf{Servir:} Dejar reposar unos minutos y servir templado.

    \end{enumerate}

\end{minipage}

\vspace{1em}

\begin{tcolorbox}[colback=orange!5, colframe=orange!50, title=Consejo, size=small]

    Queda especialmente bueno servido templado con una bola de helado de vainilla o un poco de nata montada.

\end{tcolorbox}

\end{receta}